\usepackage{bm}
%Para fotos y eso
\usepackage{graphicx}
%Para listas
\usepackage{enumitem}
%Para fórmulas y símbolos
\usepackage{amsmath}
\usepackage{amsfonts}
\usepackage{amssymb}
\usepackage{aliascnt}
%Para teoremas y demostraciones
\usepackage{amsthm}
%Para cuadraditos :)
\usepackage{mdframed}
%Headers y footers
\usepackage{fancyhdr}
%Para manejar el formato del papel
\usepackage{geometry}
%Manejo de color
\usepackage{color}
%Hipervínculos
\usepackage[colorlinks,linkcolor=blue]{hyperref}
\usepackage[nameinlink]{cleveref}
%cajitas
\usepackage[most]{tcolorbox}
\tcbuselibrary{listings, breakable, skins, theorems}
\usepackage{float}
%tikz
\usepackage{tikz,tikz-cd}

\usepackage{titlesec}
\titleformat{\chapter}[hang]{\Huge\bfseries}{Capítulo \thechapter}{1em}{}

\usepackage{graphicx,xcolor}

\usetikzlibrary{arrows.meta, positioning,calc, shapes.geometric}
%Hace setup de todos los hipervínculos del proyecto
\hypersetup{
    colorlinks=true,
    linktoc=all,
    linkcolor=black,
}
\geometry{letterpaper,margin=1in}
%SÍMBOLOS MATEMÁTICOS
%Conjuntos N,Z,Q,R,C:
\newcommand{\N}{\mathbb{N}}
    \newcommand{\Z}{\mathbb{Z}}
    \newcommand{\Q}{\mathbb{Q}}
    \newcommand{\R}{\mathbb{R}}
    \newcommand{\C}{\mathbb{C}}
    \newcommand{\F}{\mathbb{F}}
    \newcommand{\E}{\mathbb{E}}
    \newcommand{\Pp}{\mathbb{P}}
%Mathcal
\newcommand{\gL}{\mathcal{L}}
\newcommand{\gF}{\mathcal{F}}
\newcommand{\gU}{\mathcal{U}}
\newcommand{\gQ}{\mathcal{Q}}
\newcommand{\gG}{\mathcal{G}}
\newcommand{\gC}{\mathcal{C}}
\newcommand{\gT}{\mathcal{T}}
\newcommand{\gP}{\mathcal{P}}
\newcommand{\gB}{\mathcal{B}}
\newcommand{\gS}{\mathcal{S}}
\newcommand{\gO}{\mathcal{O}}
%Conjunto potencia
\newcommand{\powerset}[1]{\mathcal{P}(#1)}
%Para conjuntos
\newcommand{\set}[1]{\left\{#1\right\}}
%Función f:A->B
\newcommand{\func}[3]{#1: #2 \to #3}
%Vector cero
\newcommand{\zerov}{\overline{0}}
%Una lista indexada
\newcommand{\vectorlist}[2]{#1_1,\dots,#1_{#2}}
%Una combinación lineal
\newcommand{\lcombination}[3]{#1_1#2_1+#1_2#2_2+\dots+#1_{#3}#2_{#3}}

%Una lgualdad indexada de elementos, el tercer elemento es distinto del resto
\newcommand{\equallist}[3]{#1_1=#1_2=\dots=#1_{#2}=#3}
%Vector en F^2
\newcommand{\twovector}[2]{\begin{pmatrix}
#1\\
#2
\end{pmatrix}}
%Vector columna con n entradas
\newcommand{\coordvector}[2]{\begin{pmatrix}
#1_1\\
#1_2\\
\vdots\\
#1_{#2}
\end{pmatrix}}
%Inverso
\newcommand{\minv}[1]{#1^{-1}}
%Suma
\newcommand{\summ}[3]{\sum_{#1=#2}^{#3}}
%Representación matricial de una transformación lineal
\newcommand{\matrep}[3]{[#1]_{#2}^{#3}}
%Subsección no enumerada se añade a la tabla de contenidos.
\newcommand{\nsubsection}[1]{
\subsection*{#1}
\addcontentsline{toc}{subsection}{#1}
}

% Enteros módulo n
\newcommand{\Zn}[1]{\Z/#1\Z}
%un m-ciclo
\newcommand{\mcycle}[2]{(#1_1\spa #1_2\spa \dots \spa #1_{#2})}
%grupo cíclico
\newcommand{\cyclicgrp}[1]{\langle #1 \rangle}

\newcommand{\charsbgrp}{\ \text{char}\ }